\section{Vignettes de capacidades: el producto combina varias capacidades del modelo en un trabajo que se puede completar}

Los cinco grupos de imágenes siguientes no son una lista de funciones: muestran la \term{compositionality (composabilidad)}. Un producto útil suele invocar al mismo tiempo comprensión del lenguaje, generación de código, entorno de ejecución, renderizado visual, memoria y retroalimentación iterativa. Una capacidad que obtiene una puntuación alta en un benchmark aislado puede fallar dentro de la cadena combinada; por eso cada vignette debe leerse como un grafo de tareas de extremo a extremo.

\subsection{De la generación de código a las herramientas personales}

El ejemplo del dashboard al estilo de Stripe convierte un requisito en lenguaje natural en código de frontend y en una página interactiva. Aquí existen al menos cuatro niveles de aceptación: si se extrajeron las restricciones del requisito, si el código se ejecuta, si la jerarquía visual es legible y si, después de que el usuario hace cambios, se conserva el comportamiento existente. Si falla cualquiera de estos niveles, el producto no puede explicarse solo con la «tasa de acierto por token del código».

\lecturefigure{slide-06-stripe-dashboard.jpg}{Dashboard al estilo de Stripe generado a partir de lenguaje natural: el código y la interfaz son a la vez la salida}{00:04:41--00:05:00}

\lecturefigure{slide-07-correlation-explorer.jpg}{Explorador interactivo de correlaciones para usuarios sin formación matemática}{00:05:00--00:05:51}

El explorador de correlaciones exige además que el sistema convierta conceptos en objetos manipulables. Si el usuario arrastra la distribución de los datos pero no ve cambiar en sincronía el coeficiente de correlación, la estructura del diagrama de dispersión y los contraejemplos, la interacción es solo una «animación de demostración». Una buena eval de producto debe registrar si las variables centrales son manipulables, si la retroalimentación es inmediata, si los casos excepcionales son observables y si el usuario puede corregir, a partir de la manipulación, una idea equivocada que tenía.

\begin{importantbox}{El éxito de una tarea de extremo a extremo no es la suma de submétricas}
Sean $A_1,\dots,A_5$ los eventos de que se interprete el requisito, el código sea correcto, la ejecución tenga éxito, lo visual sea legible y el usuario complete la tarea, respectivamente. El éxito de extremo a extremo es el evento conjunto
\[
P(\text{success})=P(A_1\cap A_2\cap A_3\cap A_4\cap A_5),
\]
y no el promedio de cinco tasas de acierto. Si uno de los niveles es un umbral estricto, aunque los demás se acerquen a la puntuación máxima, el conjunto queda limitado por el eslabón más débil.
\end{importantbox}

\subsection{Imágenes, dispositivos móviles y software de costo cero}

La imagen que va de un boceto a mano a una imagen generada destaca otro tipo de interface contract: el usuario aporta una intención de baja fidelidad y el modelo completa el estilo y los detalles. Que la salida sea atractiva no significa que respete la composición, las relaciones entre objetos y la intención del autor; la evaluación debe separar aesthetic quality, instruction adherence, identity preservation y editability.

\lecturefigure{slide-08-sketch-to-image.jpg}{Del boceto a mano a la imagen generada: el modelo completa los detalles, pero el autor aún necesita controlar la dirección}{00:05:51--00:06:02}

\lecturefigure{slide-09-generative-mobile-game.jpg}{Minijuego generado al instante en el dispositivo móvil: se combinan contenido, código e interacción táctil}{00:05:52--00:06:24}

\lecturefigure{slide-10-zero-cost-software.jpg}{Demo de «la creación de software tiende a costo cero»: la interfaz móvil se genera al instante a partir de la conversación}{00:06:24--00:07:14}

Estas dos imágenes de dispositivos móviles respaldan que «el costo marginal de un prototipo disminuye», no que el software de producción ya sea gratuito. Una entrega real sigue incluyendo aclaración de requisitos, pruebas, seguridad de los datos, mantenimiento, distribución y responsabilidad. Sobre todo cuando el modelo genera software ejecutable, el entorno de ejecución debe aislarse: cuanto mayor sea la capacidad de generar código, menos se le pueden conceder al modelo por omisión permisos arbitrarios sobre la red, los archivos y las credenciales.

\teachervoice{00:04:04--00:07:36, Karina enlaza el dashboard personal, los juegos, las imágenes y el software móvil como evidencia de una «creatividad aumentada», y espera que los sistemas futuros sean más proactivos y más personalizados. Estas notas conservan su visión y añaden los límites de ingeniería: la proactividad debe diseñarse junto con los permisos, la posibilidad de revertir, el costo y la verificación.}

\begin{warningbox}{Qué es lo que más fácilmente se omite al calcular el «software de costo cero»}
El costo en tokens de generar una interfaz una vez puede ser muy bajo, pero el total cost of ownership incluye también errores en los requisitos, pruebas, infraestructura de ejecución, dependencias de terceros, auditoría de permisos, mantenimiento a largo plazo y recuperación ante incidentes. La investigación de producto debe medir el costo total de completar tareas reales, no reportar solo el precio de una generación.
\end{warningbox}

\subsection{Resumen de la sección}

La estructura común de las vignettes de capacidades es esta: la intención del usuario pasa por el modelo y la cadena de herramientas y se convierte en un artefacto observable, editable y ejecutable. Demuestran que la nueva product surface tiene valor y también exponen superficies de falla que quedan fuera de los benchmarks. La siguiente sección empieza a formular estas superficies de falla como dos scaling paradigm y dos rutas de investigación de producto.
